% ============================================================
% §5.4 Q4结果分析 — 灵敏度分析与鲁棒性检验 (v5: 题目4.1参数校正)
% 竞赛: 电工杯 B题 (2026)
% ============================================================

\section{Q4求解结果与分析：灵敏度与鲁棒性}

\subsection{三场景参数扫描设计}

按题目4.1要求"分别改变以下参数"，设计三组独立参数扰动实验，见表\ref{tab:q4_sensitivity}，每组对Q2 MILP模型独立重求解，以量化方案对人口结构、成本环境与财政预算的差异敏感性。基线为Q2确定解F+H+J，10/10全覆盖，目标值54.225。

\begin{table}[H]
\centering
\caption{三场景参数扫描灵敏度汇总（确定性求解）}
\label{tab:q4_sensitivity}
\footnotesize
\begin{tabular}{lccccc}
\toprule
\textbf{场景} & \textbf{预算(万)} & \textbf{覆盖} & \textbf{站点配置} & \textbf{平均$S$} & \textbf{年利润(万)} \\
\midrule
基线 (Q2) & 120 & 100\% & F(大)+H(中)+J(中) & 0.850 & 1~097.7 \\
\midrule
A: 人口结构冲击$^*$ & 120 & 90.0\% & B(小)+D(大)+G(小)+I(中) & 0.878 & 931.7 \\
B: 管理成本+20\%$^\dagger$ & 120 & 100\% & A(中)+F(中)+G(大) & 0.842 & 1~019.9 \\
C: 预算调整至140万 & 140 & 100\% & B(大)+D(大)+E(大) & \textbf{0.889} & 1~010.0 \\
\bottomrule
\end{tabular}

\vspace{4pt}
{\footnotesize $^*$增长率7\%$\to$8\%, $\alpha_{\text{S}\to\text{M}}$ 4.5\%$\to$5.5\%, $\alpha_{\text{M}\to\text{D}}$ 10\%$\to$9.5\% (第5年末60+人口7,953人, 失能率14.2\%). \\[2pt]
$^\dagger$仅日固定管理成本上浮20\% (建设成本不变), 对应附件3运营成本项.}
\end{table}

\subsection{场景A：人口结构冲击——四站分散化趋势}

增长率升至8\%使总人口增至7,953人，增幅5.0\%，但$\alpha_{\text{M}\to\text{D}}$从10\%降至9.5\%使失能率从14.8\%微降至14.2\%——两效应方向相反，净效应为需求从"重护理"向"重照料"倾斜。模型选择四站布局B(小)+D(大)+G(小)+I(中)，总成本113万元，覆盖率降至90\%——C小区因地理孤立未被覆盖——但满意度逆势升至0.878，增幅3.3\%，来源为小型站低利用率带来的$S_2$补偿。若固守基线方案，F站日服务量将超容量8.2\%。

\subsection{场景B：管理成本通胀——大站化集约应对}

日固定管理成本上浮20\%使运营成本权重增大，系统从"一大型两中型"切换为A(中)+F(中)+G(大)的"两中一大"布局，总成本109万元，通过规模经济——大型站日均管理成本摊薄至更大服务量——维持10/10全覆盖。满意度微降至0.842，降幅0.9\%，因G站大型化后$S_2$回升被$S_1$下降抵消——A站与G站辐射范围覆盖了部分远距离小区。

\subsection{场景C：预算释放——满意度的显著跃升}

总预算升至140万后，模型选择B(大)+D(大)+E(大)三大型站，总成本135万元，满意度从0.850跃升至0.889，增幅4.6\%，为三场景中最高。B-D-E三站沿主轴线均匀部署，$S_1$平均得分显著高于基线。20万元预算增量带来满意度增益0.039，ROI约1.95满意度点/10万元，边际报酬递减在120$\to$140万区间已开始显现。

\subsection{综合鲁棒性判定}

多维度韧性系数$\rho = 1 - |\Delta X / X_{\text{baseline}}|$，详见附录表\ref{tab:q4_resilience}、热力图\ref{fig:q4_heatmap}与雷达图\ref{fig:q4_radar}。引入三级韧性分级标准：$\rho\ge0.90$为\textbf{高韧}（冲击吸收率$\ge$90\%），$0.80\le\rho<0.90$为\textbf{中韧}（冲击吸收率80--90\%），$\rho<0.80$为\textbf{脆弱}（需预案干预）。

九组场景-维度韧性矩阵的诊断结果：覆盖率韧性0.900--1.000（全高韧），满意度韧性0.953--0.991（全高韧），利润率韧性0.849--0.929（1中韧+2高韧）。\textbf{脆弱环节识别}：场景A利润率维度$\rho=0.849$为全局最低值——人口结构冲击下利润下降15.1\%，对参数扰动最为敏感；场景A覆盖率维度$\rho=0.900$为次弱环节——C小区因地理孤立脱网。两处薄弱点均集中于场景A的人口结构扰动，提示半失能$\to$失能转化率是系统最敏感的单一参数。

\textbf{综合判定}：所有九组组合均满足$\rho>0.80$，零脆弱维度，系统整体具备强鲁棒性。韧性短板集中于利润率维度（均值0.899 vs 覆盖率0.967 vs 满意度0.970），建议在财政压力测试中将利润率作为先行预警指标。

\begin{figure}[ht]
\centering
\includegraphics[width=0.92\textwidth]{figures/figure_q4_sensitivity_combined.png}
\caption{三场景灵敏度分析全景对比。(a) 覆盖率；(b) 加权满意度；(c) 年利润。}
\label{fig:q4_sensitivity_combined}
\end{figure}

\begin{figure}[ht]
\centering
\includegraphics[width=0.62\textwidth]{figures/figure_q4_resilience_heatmap.png}
\caption{多场景多维度韧性热力图。蓝色越深韧性越强，红色越深韧性越弱。}
\label{fig:q4_heatmap}
\end{figure}

\begin{figure}[ht]
\centering
\includegraphics[width=0.52\textwidth]{figures/figure_q4_resilience_radar.png}
\caption{多维度鲁棒性雷达图。三场景各维度韧性系数均高于强鲁棒性阈值$\rho=0.80$，系统整体具备强鲁棒性。}
\label{fig:q4_radar}
\end{figure}

\subsection{对政策制定的量化启示}

\begin{enumerate}
    \item \textbf{预算140万可实现10/10全覆盖且满意度达0.889}：较基线提升4.6\%，建议将养老服务设施预算纳入财政中期规划，按建设成本指数年增3--5\%动态调整。
    \item \textbf{日管理成本+20\%不击穿全覆盖，但触发大站化}：系统通过规模经济吸收成本冲击——大型站单位服务管理成本较中型站低31\%，建议优先将大站纳入选址方案，预留应对通胀的规模弹性。
    \item \textbf{人口结构冲击触发四站分散化}：建议建立半失能$\to$失能转化率季度监测机制，转化率超11\%触发黄灯预警、超12.5\%触发红灯扩容。
\end{enumerate}

\subsection{面向全国推广的不确定因素与应对策略}

当此嵌入式养老网络从10个小区向全国铺开时，将遭遇三类题目参数空间之外的深层混沌因素。

\textbf{因素一：邻避效应。} 嵌入式站点"见缝插针"嵌入老旧小区，易引发青壮年居民对占用公共绿地、医疗废弃物、救护车噪音的群体性抗拒。应对：将养老职能"寄生"于现有社区党群中心、文体活动站，实现空间无损置换与民意柔性安抚。建模层面可将风险量化为站点社会接受度折扣因子$\gamma_i\in[0,1]$，以$C_k^{\text{build}}/\gamma_i$反映阻力成本。

\textbf{因素二：照护人力资源供应链枯竭。} Q3定价压至12.64元/次意味着服务人员单次报酬仅约8--10元，难以吸引稳定供给。建议植入"时间银行"代际接力模式：将60--70岁自理老人社区服务时间转化为可储蓄的"时间货币"，以2030年预计近5,000人自理老人池——基年4,885人按2\%年均增长外推——为劳动力蓄水池，其非货币化锚定——1小时=1积分=1元等值服务——应独立于Q3商业定价体系。

\textbf{因素三：技术颠覆的固定资产沉没风险。} 家用照护机器人一旦跨越商业量产临界点，F站45万元投资可能8--10年即面临功能性淘汰。底线思维：(1) 站点采用装配式建筑，保留拆卸重组弹性；(2) 服务设备接口标准化，确保智能终端即插即用；(3) 总投资$\ge$15\%预留为技术升级基金，使嵌入式养老网络从"一次性基建"进化为"可演进的公共服务平台"。
